\section{Paper A --- figure and table captions}

Generated by \texttt{analysis/assemble\_manuscript.py}. Numbering follows order of first citation; edit the mapping there, not here.

\subsection{Figures}

\textbf{Figure 1.} Historical observability weight \(w_t(u)\) by mechanism: absorbing loss, transient movement, and the two combined. Transient movement alone leaves \(w_t\equiv1\). First cited §4.2.

\textbf{Figure 2.} Exact cancellation is invariant to occupancy and sojourn. \(\hat\lambda/\lambda_E\) against the share of person-time spent unobservable, across sojourn lengths, under the five conditions of Corollary 2. First cited §4.2.

\textbf{Figure 3.} Analytic predictions against an independently written generative simulator. (A) Census sampling, seven population scenarios. (B) Composed with the screening and prior-testing stages. All 19 comparisons agree within \(|t|=2.20\) on 23 degrees of freedom. First cited §4.3.

\textbf{Figure S1.} Zero-bias boundary over the \((\eta_J,\eta_P)\) surface at \(q_J=3\%\), \(q_P=5\%\), with the \(r^\star=1\) contour. First cited §4.5.

\textbf{Figure S2.} Empirical recency function from the CEPHIA public-use dataset against the parametric bases. Requires data not redistributed here; see data/README.md. First cited §3.2. \emph{(not generated --- see data/README.md)}

\subsection{Tables}

\textbf{Table 1.} Recovery of the published limiting estimation error at \(w_t\equiv1\), all nine cells. First cited §4.1.

\textbf{Table 2.} Attenuation and the composed zero-bias boundary against absorbing loss \(\mu\). First cited §4.4.

\textbf{Table 3.} Cancellation under stationary frailty mixtures of increasing skew in movement propensity. First cited §4.2.

\textbf{Table S1.} CEPHIA MDRI by algorithm and subtype. Requires data not redistributed here. First cited §3.2. \emph{(not generated --- see data/README.md)}

\textbf{Table S2.} Covariate reweighting against the duration mechanism, three target-population cases. First cited §4.8.

\textbf{Table S3a.} Zero-bias boundary against a common relative acquisition hazard \(\eta\). First cited §4.5.

\textbf{Table S3b.} Break-even \(\eta\) by site county. Illustrative; not an epidemiologic claim about any county. First cited §4.7.

\textbf{Table S4.} Boundary invariance across recency bases, by inter-test process. First cited §4.6.
